% !TEX root = chapter-standalone.tex


\section{Semicategories}
\label{sec:semicategories}

Semicategories are  ``categories which don't necessarily have \SY{identity morphisms}''.
A good analogy is that \SY{semicategories} are to categories as  \SY{semigroups} are to \SY{monoids}..
A \SY{monoid} is a \SY{semigroup} that additionally has an identity element, and similarly a category is a \SY{semicategory} that additionally has \SY{identity morphisms}.

\begin{ctdefinition}[Semicategory]
    \label{def:semicategory-compact}
    A \maindef{semicategory}~\CatC is specified by:
    \begin{body}
        \constit
        \begin{enumerate}
            \item \label{semicategory-objects}
                  Objects: A collection\footnotemark~$\ObC$ whose elements are called \emph{objects}.
            \item \label{def:hom-set}
                  Morphisms: For every pair of objects~$\Obja,\Objb$ in $\ObC$, there is a set called a ``hom-set'' and indicated as~$\HomSet{\CatC}{\Obja}{\Objb}$, elements of which are called \emph{morphisms} and denoted $\mora \colon \Obja \mto \Objb$.

                  For such an $\mora$, we call $\Obja$ its \emph{source} and $\Objb$ its \emph{target}.
            \item \label{semicategory-composition}
                  Composition operations: For every three objects~$\Obja,\Objb,\Objc$ in~$\ObC$ there is a composition map
                  \begin{equation}\label{eq:semicategory-composition}
                      \mthenof{\Obja,\Objb,\Objc}\colon \HomSet{\CatC}{\Obja}{\Objb} \cartprod \HomSet{\CatC}{\Objb}{\Objc}\mto \HomSet{\CatC}{\Obja}{\Objc}.
                  \end{equation}
                  We usually just write~$\mthen$ instead of~$\mthenof{\Obja,\Objb,\Objc}$:
                  \begin{equation}
                      \prfperiod{
                          \mora \colon \Obja \mto \Objb
                      }{
                          \morb\colon \Objb \mto \Objc
                      }{
                          (\morab )\colon \Obja \mto \Objc
                      }
                  \end{equation}
                  The morphism~$\morab$ is called the \emph{composition} of~$\mora$ and~$\morb$.
                  %Given any morphism~${\mora} \setin  \HomSet{\CatC}{\Obja}{\Objb}$ and any morphism~${\morb} \setin \HomSet{\CatC}{\Objb}{\Objc}$, there exists a morphism~$\mora \mthen \morb\setin \HomSet{\CatC}{\Obja}{\Objc}$ which is the \emph{composition of~$\mora$ and~$\morb$}.
        \end{enumerate}
        \condit
        \begin{enumerate}
            \item \label{semicategory-associativity}
                  Associativity: it holds that
                  % for any morphisms~$\mora\setin \HomSet\CatC\Obja\Objb$,~$\morb\setin \HomSet\CatC\Objb\Objc$, and~$\morc\setin \HomSet\CatC\Objc\Objd$,
                  \begin{equation}
                      \prfperiod{
                          \mora\colon \Obja \mto \Objb
                      }{
                          \morb\colon \Objb \mto \Objc
                      }{
                          \morc\colon \Objc \mto \Objd
                      }{
                          (\morab)
                          \mthen \morc= \mora \mthen (\morb \mthen \morc)
                      }
                  \end{equation}
        \end{enumerate}
    \end{body}
\end{ctdefinition}

\subsection{First examples}



%Now let's view electrical chords as forming a semicategory. For simplicity (and for the ``C''), let's call our semicategory-to-be by the name $\CatC$, same as in the definition above. To define $\CatC$, let the collection $\ObC$ of objects to be the collection of all electrical interface types: one object for each type. Next, given any two types $\Obja$ and $\Objb$ (\text{e.g.}, $\Obja = \TypeIrish$ and $\Objb = \TypeSwiss$), let's define the set of morphisms $\HomSet{\CatC}{\Obja}{\Objb}$ to be the set of all electrical chords that have the type $\Obja$ as their incoming interface and type $\Objb$ as their outgoing interface.

%\let\oldmora\mora
%\renewcommand{\mora}{{\colTransmuter \oldmora}}
%\let\oldmorb\morb
%\renewcommand{\morb}{{\colTransmuter \oldmorb}}
%\let\oldmorc\morc
%\renewcommand{\morc}{{\colTransmuter \oldmorc}}
%\let\oldObja\Obja
%\renewcommand{\Obja}{{\colTransmuted \oldObja}}
%\let\oldObjb\Objb
%\renewcommand{\Objb}{{\colTransmuted \oldObjb}}
%\let\oldObjc\Objc
%\renewcommand{\Objc}{{\colTransmuted \oldObjc}}
%\let\oldObjd\Objd
%\renewcommand{\Objd}{{\colTransmuted \oldObjd}}
%
%\let\oldOb\Ob
%\renewcommand{\Ob}{{\colTransmuted \oldOb}}
%
%\let\oldHom\Hom
%\renewcommand{\Hom}{{\colTransmuter \oldHom}}



    \begin{example}
    \label{exa:semicat-plant-states}
    This example is an extension of \cref{exa:plant-trafo-semigroup}.
    We will describe a \SY{semicategory}~\CatC with two objects,~$\ObC = \makeset{ \setA, \setB }$.
    Each of the objects is a set which describes possible states of a plant.
    Let
    \begin{equation}
        \setA = \makeset{ \sprout, \yng, \mature, \old, \dead },
    \end{equation}
    and
    \begin{equation}
        \setB = \makeset{ \alive, \dead }.
    \end{equation}
    Also, let~$\mapa \colon \setA \sto \setA$ be the function with
    \begin{equation}
        \begin{aligned}
            \mapa(\sprout) & = \yng, \\
            \mapa(\yng)    & = \mature, \\
            \mapa(\mature) & = \old, \\
            \mapa( \old)   & = \dead, \\
            \mapa(\dead)   & = \dead;
        \end{aligned}
    \end{equation}
    let~$\morb \colon \setB \sto \setB$ be the function with
    \begin{equation}
        \begin{aligned}
            \mapb(\alive) & = \dead, \\
            \mapb(\dead)  & = \dead;
        \end{aligned}
    \end{equation}
    and let~$\morc \colon \setA \sto \setB$ be the function with
    \begin{equation}
        \begin{aligned}
            \mapc(\sprout) & = \alive, \\
            \mapc(\yng)    & = \alive, \\
            \mapc(\mature) & = \alive, \\
            \mapc( \old)   & = \alive, \\
            \mapc (\dead)  & = \dead.
        \end{aligned}
    \end{equation}
    For our sets of morphisms, we let
    \begin{equation}
        \begin{aligned}
            \HomSet{\CatC}{\setA}{\setA} & = \makeset{\mora, \mora^2, \mora^3, \mora^4 }, \\
            \HomSet{\CatC}{\setB}{\setB} & = \makeset{ \morb }, \\
            \HomSet{\CatC}{\setA}{\setB} & = \makeset{ \morc,\ \mora \mthen \morc,\ \mora^2 \mthen \morc, \ \mora^3 \mthen \morc,
            \ \mora^4 \mthen \morc}, \\
            \HomSet{\CatC}{\setB}{\setA} & = \Emptyset;
        \end{aligned}
    \end{equation}
    and for the composition operations, we define these to be the usual composition of functions, which we know obeys the \SY{associative law}.
    Thus, we have a \SY{semicategory}.
    \todojira{269}{\bernina: Finish writing the above example involving plant states by adding a nice figure to illustrate it}
\end{example}

\todotext{\bernina: @JL: Seems like it would be good to add more examples here}

\todotextjira{394}{\bernina: @JL: Add remark about relationship between \SY{semicategories} and semigroups}

\todotext{J: Maybe make an example here where morphisms are taxi trips and their costs, calculated as an initial fee, plus a rate for time and kilometers each.
    There are no identity morphisms because any trip incurs a non-zero cost.
}

